% =====================================================================
% APPENDIX C: Addressing Complex Engineering Activities
% (RESTRUCTURE_PLAN_V2, appx-kp owner)
% Ported from old mainmatter/ch5.tex Section 5.3 (Addressing Complex
% Engineering Activities), with old Section 5.4's closing sentences on
% the A attributes folded into this appendix's closing paragraph.
% Chapter cross-references updated from old hard-coded numbers to the
% new 9-chapter labels; the two mentions of the old socio-economic
% chapter are phrased descriptively since that chapter's label had not
% been assigned yet.
% =====================================================================
\chapter{Addressing Complex Engineering Activities}
\label{app:cea}

The Washington Accord also characterizes complex engineering
activities by five A attributes, distinct from the P attributes of a
complex problem discussed in Appendix~\ref{app:cep}. An activity is
complex when carrying it out demands diverse resources, conflicting
requirements, innovation, and the other characteristics set out below.
The work we carried out in this thesis exhibits each of the five.

\textbf{A1, range of resources.} We used several distinct kinds of
resources: rented cloud GPUs for training and ablation, multiple
public stereo datasets for training and cross-domain evaluation, a
physical calibrated stereo camera rig for real-data experiments, an
embedded accelerator as the deployment target, and the open-source
deep-learning tooling on which the implementation of
Chapter~\ref{ch:implementation} is built. The management of these
resources, including their cost, is documented in this thesis's
management and finance account.

\textbf{A2, level of interaction.} The project's central technical
interactions are the conflicts among accuracy, latency, and memory:
improving one routinely degrades another. We could not resolve these
conflicts analytically. We resolved them instead through the
controlled ablation methodology of Chapter~\ref{ch:design}, under
which every candidate change was measured on all metrics jointly and
adopted only when the trade-off was favorable. Chapter~\ref{ch:results}
records these studies.

\textbf{A3, innovation.} The activity was creative rather than
routine. The architecture of Chapter~\ref{ch:design} is a novel
composition of building blocks adapted from distinct families of the
literature. It also includes the fail-soft fusion of the geometry
encoding volume, a mechanism we designed in this work so that the
added evidence sharpens the estimate when it is reliable and degrades
gracefully, rather than catastrophically, when it is not.

\textbf{A4, consequences to society and the environment.} The activity
has identifiable consequences beyond the laboratory: on-device
privacy, robot and drone safety, eye-safe passive sensing,
accessibility of low-cost depth, energy saved by avoiding cloud round
trips, and a transparent training footprint. This thesis's
socio-economic impact analysis examines these consequences.

\textbf{A5, familiarity.} The activity extends beyond textbook
material and prior experience. Its subject matter, compact recurrent
stereo architectures, cross-domain generalization, and edge-inference
optimization, exists only in the recent research literature we
reviewed in Chapter~\ref{ch:lit}. Engaging with it required reading,
evaluating, and adapting published research rather than applying
established practice.

This appendix has shown that the work carried out in this thesis
constitutes complex engineering activities in the Washington Accord
sense. The five A attributes are exhibited through the range of
resources we used, the measured resolution of conflicting
requirements, the novel block composition with fail-soft fusion, the
societal and environmental consequences examined in the
socio-economic impact analysis, and the beyond-textbook research
grounding of Chapter~\ref{ch:lit}. Together with Appendix~\ref{app:cep},
this completes the case that the thesis addresses both a complex
engineering problem and complex engineering activities.

% SOURCES: old mainmatter/ch5.tex Section 5.3 (Addressing Complex
% Engineering Activities, all five A-attribute paragraphs carried over
% with judgment) and part of Section 5.4 (Summary, the A-attribute
% closing sentences folded into this appendix's closing paragraph).
% DROPPED: the "Section 5.2 does X, Section 5.3 does Y" roadmap
% sentence from old Section 5.1 was dropped as redundant once the
% material was split across three separate appendix files.
% MOVED-AWAY: the CEP material (old Section 5.2) is in app_b.tex, and
% the knowledge-profile argument is in app_a.tex, both owned by this
% same agent (appx-kp).
